\section{Description}

The task is to implement the radix sort algorithm for \enquote{date--value} pairs.

The main idea of radix sort is to split every input element into digits and to apply
a stable counting sort to each digit, from the least significant digit to the most
significant one. After all digits have been processed, the array is fully sorted.
The complexity of the algorithm is $O(n \cdot k)$, where $k$ is the number of digits
and $n$ is the number of elements.

\pagebreak

\section{Source code}

\subsection{Auxiliary structures}

The key is a date in the format \texttt{DD.MM.YYYY}, which can be uniquely represented
as an integer by the formula $year \times 10000 + month \times 100 + day$. This encoding
preserves the order of dates under numeric comparison.

Since the number of input lines is not known in advance, a dynamic array
\texttt{Vector<T>} that doubles its capacity on overflow is implemented. It follows
the rule of five: destructor, copy constructor, copy assignment operator,
move constructor and move assignment operator. The structure \texttt{Pair} is a plain
aggregate of a \texttt{std::string} and a \texttt{std::uint32\_t} and follows the rule of zero:
all special member functions are generated by the compiler.

Input is read line by line. Blank lines and lines consisting only of whitespace are skipped,
and a trailing carriage return is stripped. A line without a tab separator, a value longer
than 64 characters or an invalid date is rejected: the program prints a message starting with
\texttt{ERROR:} to the standard error stream and exits with code 1.

\begin{table}[!ht]
\centering
\begin{tabular}{|P{7.5cm}|P{7.5cm}|}
\hline
\rowcolor{lightgray}
\multicolumn{2}{|c|}{include/vector.hpp} \\
\hline
\texttt{template <typename T> class Vector} & Dynamic array with the rule of five; methods \texttt{pushBack}, \texttt{size}, \texttt{capacity}, \texttt{empty}, \texttt{data}, \texttt{begin}, \texttt{end}, \texttt{swap}, \texttt{operator[]} \\
\hline
\rowcolor{lightgray}
\multicolumn{2}{|c|}{include/pair.hpp, src/pair.cpp} \\
\hline
\texttt{struct Pair} & Key-value pair: stores the original line \texttt{raw} and the numeric key \texttt{key} of type \texttt{std::uint32\_t} \\
\hline
\texttt{bool isBlank(std::string\_view)} & Checks whether a line consists only of whitespace characters \\
\hline
\texttt{std::uint32\_t parseDateKey(std::string\_view)} & Parses a date \texttt{DD.MM.YYYY} and encodes it as an integer key \\
\hline
\texttt{Pair parsePair(std::string)} & Parses one input line into a pair \\
\hline
\texttt{Vector<Pair> readPairs(std::istream\&)} & Reads all non-blank lines of the input \\
\hline
\texttt{void writePairs(std::ostream\&, const Vector<Pair>\&)} & Prints the pairs in the original form \\
\hline
\rowcolor{lightgray}
\multicolumn{2}{|c|}{include/radix\_sort.hpp, src/radix\_sort.cpp} \\
\hline
\texttt{void countingPass(Vector<Pair>\&, Vector<Pair>\&, int)} & One counting sort pass over the given byte of the key \\
\hline
\texttt{void radixSort(Vector<Pair>\&)} & Radix sort: four passes over the bytes of \texttt{std::uint32\_t} \\
\hline
\end{tabular}
\end{table}

\lstinputlisting[caption={include/vector.hpp}]{../include/vector.hpp}

\lstinputlisting[caption={include/pair.hpp}]{../include/pair.hpp}

\lstinputlisting[caption={src/pair.cpp}]{../src/pair.cpp}

\subsection{Counting sort over one byte}

The function \texttt{countingPass} performs one stable pass over the given byte of the key.
First, an array of frequencies is built (256 elements, one for each possible byte value),
then it is transformed into an array of prefix sums. The elements are placed into the output
array in reverse order, which guarantees stability.

\subsection{Radix sort}

The function \texttt{radixSort} applies \texttt{countingPass} successively to bytes 0, 1, 2, 3
(from the least significant to the most significant). To avoid repeated allocations, a
preallocated \texttt{buffer} is used: after each pass the input array and the buffer are
exchanged with \texttt{swap}, so no elements are copied back.

\lstinputlisting[caption={include/radix\_sort.hpp}]{../include/radix_sort.hpp}

\lstinputlisting[caption={src/radix\_sort.cpp}]{../src/radix_sort.cpp}

\subsection{Entry point}

The function \texttt{main} reads the pairs, sorts them with \texttt{radixSort} and prints the
result. Any parsing error is reported to the standard error stream.

\needspace{18\baselineskip}
\lstinputlisting[caption={app/main.cpp}]{../app/main.cpp}

\pagebreak

\section{Console}

\begin{alltt}
# make build PRESET=release
# build/release/lab1/lab1_main
01.01.2000    Alice
15.06.1995    Bob
03.03.2010    Charlie
20.11.1988    Diana
^D
20.11.1988    Diana
15.06.1995    Bob
01.01.2000    Alice
03.03.2010    Charlie
\end{alltt}

\pagebreak
